%%%%%%%%%%%%%%%%%%%%%%preface.tex%%%%%%%%%%%%%%%%%%%%%%%%%%%%%%%%%%%%%%%%%
% 
% Preface for The Trustworthy AI Architect
%
%%%%%%%%%%%%%%%%%%%%%%%% Springer Nature %%%%%%%%%%%%%%%%%%%%%%%%

\preface

The transition from a set of occasional blog posts into a cohesive, 25-chapter monograph began with a simple observation: while the world is racing to deploy Artificial Intelligence at scale, the engineering discipline for making that AI \textit{trustworthy} is still in its infancy. This book is the result of a journey to bridge that gap.

\textit{The Trustworthy AI Architect} was born from the "Building Trustworthy AI" series, a collection of reflections on privacy, security, and verification originally written for a global and Vietnamese audience. Over the course of these chapters, my goal is to move beyond the high-level policy debates and dive into the specific architectural patterns---from Differential Privacy to Trusted Execution Environments---that make trust a verifiable system property rather than just a marketing promise.

This book is intended for practitioners, students, and researchers who believe that accuracy is no longer the sole metric of success. It is for those who realize that a model is only as valuable as the privacy it protects and the safety it guarantees.

As we deploy these systems across Southeast Asia and the world, it is my hope that this handbook serves as a guide for the next generation of architects who will build the resilient, private, and secure foundations of our future.

\vspace{\baselineskip}
\begin{flushright}\noindent
Ho Chi Minh City, Vietnam \hfill {\it Phat T. Tran-Truong (Leo)}\\
April 2024 \hfill {\it }\\
\end{flushright}
